\chapter{Curve Construction}\label{ch:rc:curve-construction}

On a Friday evening a rates desk replaces the interpolation of its dollar
curve: zero rates had been joined by straight lines, and from Monday they will
be joined by the method the US Treasury itself adopted in 2021. Every input
quote on Monday morning is exactly where it was on Friday, every input
instrument still reprices to the last decimal, and yet the desk's book of
seven \omterm{def:rc:curve-construction:pillar}{off-pillar} swaps, five billion dollars of notional, shows a loss of
90 thousand dollars, and the hedge the risk system asks for on the fifteen-year
swap has grown by two thirds. Nothing in the market moved; the curve between
the quotes did. A curve is a model, fitted to a handful of prices and used to
value everything between them. This chapter builds one from the instruments
up, and measures what each choice costs in value and in hedges. One Quant Book
2, chapter 9, bootstrapped a single overnight-swap curve with one
interpolation; here the curve takes deposits, futures and swaps together, is
solved globally, and is interpolated in four ways.

\section{Choosing the instruments}

\begin{definition}[Pillar, instantaneous forward rate]\label{def:rc:curve-construction:pillar}
A \emph{pillar}\index{pillar} of a curve is a date at which the curve carries a
free parameter, here the continuously compounded zero rate $z(T)$, so that
$P(T)=e^{-z(T)T}$; the curve between pillars is fixed by an interpolation rule.
The \emph{instantaneous forward rate}\index{instantaneous forward rate} seen at
$t$ for time $T$ is
\[
f(t,T) = -\frac{\partial}{\partial T}\ln P(t,T),
\qquad
P(t,T)=\exp\Bigl(-\int_t^T f(t,u)\,du\Bigr),
\]
the rate at which one can lock in, at $t$, borrowing over $[T, T+dT]$.
\end{definition}

\begin{definition}[Curve calibration]\label{def:rc:curve-construction:calibration}
\emph{Curve calibration}\index{curve calibration} is the solution for the
\omterm{def:rc:curve-construction:pillar}{pillar} parameters of a curve such that a chosen set of instruments, each
priced on the curve, reproduces its market quote. With one \omterm{def:rc:curve-construction:pillar}{pillar} per
instrument it is a square system of equations; the general notion of
calibration, fitting a model's parameters to prices, is treated in One Quant
Book 4, chapter 24.
\end{definition}

The instruments are chosen for liquidity and so that each maturity is covered
once. At the front of a dollar curve sit overnight fixings and short
overnight-index deposits; then three-month SOFR futures, whose reference
quarters tile the first one or two years; then par OIS swaps
(One Quant Book 2, chapters 8 and 9). The rule is \emph{no overlaps}: if four
futures already fix the forwards to December 2027, a one-year swap maturing in
September 2027 adds a second, slightly different price for a period the curve
has already priced. Two quotes for one period cannot both be matched unless
they agree; the builder either drops one or fits both in a least-squares sense
and gives up exact repricing. \Cref{fig:rc:curve-construction:instruments}
shows the chapter's choice.

\begin{omfigure}
\begin{tikzpicture}[x=4.4cm,y=1cm,font=\small,
  bar/.style={draw=omDef,fill=omDef!15,thick,minimum height=0.36cm},
  lab/.style={font=\footnotesize,anchor=west}]
\draw[->,thick] (0,-1.2) -- (2.25,-1.2) node[right] {years};
\foreach \x/\l in {0/0,0.5/0.5,1/1,1.5/1.5,2/2} {\draw (\x,-1.28) -- (\x,-1.12); \node[below,font=\footnotesize] at (\x,-1.28) {\l};}
% deposits
\draw[bar] (0,2.55) rectangle (0.082,2.85); \node[lab] at (0.10,2.70) {1M deposit};
\draw[bar] (0,2.05) rectangle (0.249,2.35); \node[lab] at (0.27,2.20) {3M deposit};
% futures
\foreach \a/\b/\n in {0.2137/0.463/F1,0.463/0.7123/F2,0.7123/0.9616/F3,0.9616/1.211/F4}
  {\draw[draw=omThm,fill=omThm!15,thick] (\a,1.05) rectangle (\b,1.45); \node[font=\footnotesize] at ({(\a+\b)/2},1.25) {\n};}
\node[lab] at (1.24,1.25) {SOFR futures (IMM quarters)};
% swaps
\draw[draw=omProp,fill=omProp!15,thick] (0,0.05) rectangle (2.0027,0.45); \node[font=\footnotesize] at (1.0,0.25) {2Y par OIS swap};
\node[lab,align=left] at (2.03,0.25) {then 3Y, 4Y, 5Y, 7Y,\\10Y, \dots, 30Y};
\draw[dashed,gray,thick] (0,-0.75) rectangle (1.0,-0.45); \node[lab,gray] at (1.03,-0.6) {1Y swap: not used, overlaps F1--F4};
% pillars
\foreach \x in {0.082,0.249,0.463,0.7123,0.9616,1.211,2.0027} {\fill[omExo] (\x,-1.2) circle (1.6pt);}
\end{tikzpicture}

\omcaption{The instruments of the chapter's dollar curve over its first two
years, and their \omterm{def:rc:curve-construction:pillar}{pillars} (dots on the axis, one at each instrument's
maturity). The futures start on the third Wednesday of December 2026 and tile
five quarters; the one-year swap would price a period the futures already
price, and is left out.}\label{fig:rc:curve-construction:instruments}
\end{omfigure}

\begin{dated}{2026-09}{Published curves}\label{dat:rc:curve-construction:published}
The US Treasury's official daily par yield curve has been built with a
monotone convex method since 6 December 2021, replacing a quasi-cubic Hermite
spline; its inputs are indicative bid-side quotes for the most recently
auctioned bills, notes and bonds, collected by the New York Fed near 15:30.
The ECB publishes two euro-area government curves, from AAA-rated bonds and
from all central-government bonds, every TARGET business day at 12:00 CET,
each as the six parameters of a Svensson model
(\cref{def:rc:curve-construction:ns}) and the rates they imply.
\end{dated}

\section{Calibration: curve bootstrapping and the global fit}

Each instrument's quote is a function of the discount factors up to its
maturity. A deposit from spot to $T$ at simple rate $R$ (ACT/360) requires
$P(T) = 1/(1+R\delta)$; a future on the quarter $[T_1,T_2]$, once its
convexity adjustment is removed (One Quant Book 2, chapter 8), fixes
$P(T_1)/P(T_2) = 1+F\delta$; a par swap fixes $P(T_0)-P(T_n)=K\,\sum_i
\delta_iP(t_i)$. With $n$ instruments and $n$ \omterm{def:rc:curve-construction:pillar}{pillars} the system
$m_i(z_1,\dots,z_n)=q_i$ is square.

\begin{proposition}[When curve bootstrapping works]\label{prop:rc:curve-construction:triangular}
If the interpolation is \emph{local}, in the sense that the curve on
$[T_{j-1},T_j]$ depends only on the \omterm{def:rc:curve-construction:pillar}{pillars} $j-1$ and $j$, and each instrument
$i$ has maturity $T_i$, then $\partial m_i/\partial z_j = 0$ for $j>i$: the
Jacobian is lower triangular, and the system can be solved one \omterm{def:rc:curve-construction:pillar}{pillar} at a
time, each a one-dimensional root search. This is curve bootstrapping. Linear
interpolation of zero rates and of log discount factors are local; a cubic
spline is not, and the monotone convex method is local except through the
forwards it estimates at the two neighbouring \omterm{def:rc:curve-construction:pillar}{pillars}.
\end{proposition}
\begin{proof}
Instrument $i$ only reads discount factors at dates up to $T_i$; under a local
rule these depend only on \omterm{def:rc:curve-construction:pillar}{pillars} $1,\dots,i$.
\end{proof}

\begin{method}[Global curve calibration by Newton's method]\label{met:rc:curve-construction:newton}
Start from $z^{(0)}_j=q_j$. At each step compute the residuals $r_i = m_i(z)
- q_i$ and the Jacobian $J_{ij}=\partial m_i/\partial z_j$ by finite
differences; update $z \leftarrow z - J^{-1}r$; stop when $\max_i|r_i|<10^{-12}$.
The same code serves every interpolation, local or not, and the final Jacobian
is kept: it converts sensitivities to \omterm{def:rc:curve-construction:pillar}{pillar} rates into sensitivities to
quotes (\cref{ch:rc:rates-risk}).
\end{method}

\omcode{code/firm/curvebuild/firm_curvebuild.py}{256}{284}{Global calibration of all pillar zero rates by Newton's method.}\label{lst:rc:curve-construction:calibrate}

On the seventeen instruments of \cref{fig:rc:curve-construction:instruments}
Newton's method converges in three iterations for all four interpolations,
to residuals below $10^{-15}$. Speed is not the point here; what the global
form buys is that the interpolation can be changed without changing the solver.

\section{Interpolation and the forward curve}

Every interpolation reprices the inputs; they differ in what they say between
\omterm{def:rc:curve-construction:pillar}{pillars}, and the forward curve is where the differences show.

\begin{proposition}[Forwards from zero rates]\label{prop:rc:curve-construction:fwd}
With $y(T) = z(T)T = -\ln P(T)$, the instantaneous forward is $f(0,T) =
y'(T) = z(T) + T z'(T)$. The average forward over a \omterm{def:rc:curve-construction:pillar}{pillar} interval is fixed by
the two \omterm{def:rc:curve-construction:pillar}{pillars} alone: $\bar f_j = (y_j - y_{j-1})/(T_j-T_{j-1})$, the
\emph{discrete forward}.
\end{proposition}

\begin{definition}[Flat-forward interpolation]\label{def:rc:curve-construction:flatfwd}
\emph{Flat-forward interpolation}\index{flat-forward interpolation} makes
$\ln P(T)$ linear between \omterm{def:rc:curve-construction:pillar}{pillars}: the instantaneous forward is constant on
each interval and equal to its discrete forward, and jumps at the \omterm{def:rc:curve-construction:pillar}{pillars}.
\end{definition}

Linear interpolation of zero rates looks smoother on a zero-rate chart and is
worse on a forward chart: by \cref{prop:rc:curve-construction:fwd} the forward
is $z + Tz'$, and $z'$ jumps at each \omterm{def:rc:curve-construction:pillar}{pillar}, multiplied by $T$. At the
chapter's two-year \omterm{def:rc:curve-construction:pillar}{pillar} ($z = 3.2883\%$) the zero curve falls on its left
and rises on its right; the forward jumps from $2.964\%$ to $3.331\%$, a
37-basis-point step made by the rule, not by the market. A natural cubic
spline through the zero rates (One Quant Book 4, chapter 28) removes the
jumps but, being a global fit, lets a change at one \omterm{def:rc:curve-construction:pillar}{pillar} ripple through
all of them.

\begin{definition}[Monotone convex interpolation]\label{def:rc:curve-construction:monotone}
\emph{Monotone convex interpolation}\index{monotone convex interpolation}
(Hagan and West, 2006) builds the forward curve directly. On each interval
it keeps the discrete forward $\bar f_j$, estimates the instantaneous forward
$f_j$ at each \omterm{def:rc:curve-construction:pillar}{pillar} as the time-weighted average of the two adjacent discrete
forwards, and fills the interval with $f(T)=\bar f_j + G(x)$,
$x=(T-T_{j-1})/(T_j-T_{j-1})$, where $G$ is a quadratic, or a quadratic joined
to a constant, with $G(0)=f_{j-1}-\bar f_j$, $G(1)=f_j-\bar f_j$ and
$\int_0^1G=0$, chosen from four cases so that $G$ never overshoots its end
values. The forward curve is continuous, reproduces every discrete forward,
and stays within the range of the neighbouring forwards.
\end{definition}

\omcode{code/firm/curvebuild/firm_curvebuild.py}{123}{138}{The four cases of the monotone convex interpolant on one interval.}\label{lst:rc:curve-construction:mc}

\begin{omfigure}
\begin{tikzpicture}
\begin{axis}[width=12cm,height=6cm,
  xlabel={maturity $T$ (years)},ylabel={$f(0,T)$ (\%)},
  tick label style={font=\small},label style={font=\small},
  xmin=0,xmax=30,ymin=2.9,ymax=4.75,grid=major,grid style={gray!25},
  legend columns=4,legend style={font=\footnotesize,at={(0.5,-0.25)},anchor=north,draw=none,fill=none,/tikz/every even column/.append style={column sep=8pt}},
  table/col sep=comma]
\addplot[gray,thick] table[x=t,y=lz]{figdata/rates-credit-risk/01-curve-construction/forwards.csv};
\addplot[omDef,thick] table[x=t,y=ff]{figdata/rates-credit-risk/01-curve-construction/forwards.csv};
\addplot[omProp,thick,dashed] table[x=t,y=cs]{figdata/rates-credit-risk/01-curve-construction/forwards.csv};
\addplot[omThm,very thick] table[x=t,y=mc]{figdata/rates-credit-risk/01-curve-construction/forwards.csv};
\legend{linear zero,flat forward,cubic on zeros,monotone convex}
\end{axis}
\end{tikzpicture}

\omcaption{Four instantaneous forward curves from the same seventeen quotes;
every one reprices every input exactly. Linear zero rates give a sawtooth;
flat forwards give steps; the cubic spline and the monotone convex curve are
continuous and nearly coincide at this scale: they differ in how they respond
to a change of one quote (\cref{fig:rc:curve-construction:locality}).
Data: the chapter's illustrative dollar curve and
tutorial.}\label{fig:rc:curve-construction:forwards}
\end{omfigure}

For government bonds, which are many, noisy and not all at par, a smooth
parametric curve fitted by least squares is usual instead of one \omterm{def:rc:curve-construction:pillar}{pillar} per
instrument.

\begin{definition}[Nelson--Siegel and Svensson curves]\label{def:rc:curve-construction:ns}
A \emph{Nelson--Siegel curve}\index{Nelson--Siegel curve} is the zero curve
\[
z(T) = \beta_0 + \beta_1\,\frac{1-e^{-T/\tau_1}}{T/\tau_1}
 + \beta_2\Bigl(\frac{1-e^{-T/\tau_1}}{T/\tau_1}-e^{-T/\tau_1}\Bigr),
\]
with long-run level $\beta_0$, short end $\beta_0+\beta_1$ and a hump of size
$\beta_2$ near $\tau_1$; its forward curve is $f(0,T) = \beta_0 + \beta_1
e^{-T/\tau_1} + \beta_2 (T/\tau_1)e^{-T/\tau_1}$. Svensson's extension adds a
second hump term $\beta_3$ with its own $\tau_2$.
\end{definition}

\begin{example}[The ECB's AAA curve]\label{ex:rc:curve-construction:ecb}
On 22 September 2026 the ECB's AAA Svensson parameters were $\beta_0=1.6306$,
$\beta_1=0.6958$, $\beta_2=2.5849$, $\beta_3=6.6622$ (per cent), $\tau_1=1.0632$,
$\tau_2=14.4456$ years. They give a ten-year zero rate of $3.4527\%$, which is
the ECB's published ten-year spot rate, and the one-, five- and thirty-year
published rates to six decimals; the instantaneous forward rises from
$2.33\%$ at the short end to a peak of $4.08\%$ near fourteen years
(\cref{fig:rc:curve-construction:svensson}).
\end{example}

\begin{omfigure}
\begin{tikzpicture}
\begin{axis}[width=11cm,height=5.4cm,
  xlabel={maturity $T$ (years)},ylabel={rate (\%)},
  tick label style={font=\small},label style={font=\small},
  xmin=0,xmax=30,ymin=2,ymax=4.4,grid=major,grid style={gray!25},
  legend columns=3,legend style={font=\footnotesize,at={(0.5,-0.25)},anchor=north,draw=none,fill=none,/tikz/every even column/.append style={column sep=8pt}},
  table/col sep=comma]
\addplot[omDef,very thick] table[x=t,y=zero]{figdata/rates-credit-risk/01-curve-construction/svensson.csv};
\addplot[omThm,thick,dashed] table[x=t,y=fwd]{figdata/rates-credit-risk/01-curve-construction/svensson.csv};
\addplot[only marks,mark=*,mark size=2.2pt,omExo] table[x=t,y=zero]{figdata/rates-credit-risk/01-curve-construction/ecbpoints.csv};
\legend{zero rate,instantaneous forward,ECB published spot rates}
\end{axis}
\end{tikzpicture}

\omcaption{The euro-area AAA government curve of 22 September 2026, from the
six Svensson parameters the ECB publishes; the dots are the spot rates the ECB
publishes for one, five, ten and thirty years, which the formula reproduces.
Source: ECB statistics (ECB Data Portal, series
YC.B.U2.EUR.4F.G\_N\_A.SV\_C\_YM).}\label{fig:rc:curve-construction:svensson}
\end{omfigure}

\section{The front end: futures, turns and meeting dates}

The first two years of a dollar curve are built from instruments that price
forward periods, not rates from spot. A future on an IMM quarter fixes one
quarter's forward rate, after its convexity adjustment; four of them fix four
consecutive discrete forwards, and the interpolation only decides the shape
inside each quarter. Two features of the short end are not smooth and should
not be interpolated as if they were. The turn
of One Quant Book 2, chapter 2, is a jump in the overnight rate over a
year-end or quarter-end: a curve builder that knows it adds an extra \omterm{def:rc:curve-construction:pillar}{pillar}
around the turn date, or a separate spread over those days, so that the one
expensive night is not smeared over a quarter. The implied policy path
is a step function: the overnight rate moves at central-bank meetings and not
between them, so front-end curves are often built with \omterm{def:rc:curve-construction:flatfwd}{flat-forward
interpolation} between meeting dates, one \omterm{def:rc:curve-construction:pillar}{pillar} per meeting (Book 2's build
\texttt{firm.meetings}). Smooth interpolation is right where the market has no
calendar, from two or three years on.

\begin{remark}[Which curve for which job]\label{rem:rc:curve-construction:which}
Flat forwards between meeting dates at the front, a smooth rule such as
monotone convex beyond, and a parametric curve for relative value in
government bonds (a bond rich or cheap to the fitted curve) are complementary,
not competing, choices. Any of them fails if an instrument is included twice
for one period or a quote is stale.
\end{remark}

\section{Hedging consequences of an interpolation}

\begin{definition}[Interpolation locality]\label{def:rc:curve-construction:locality}
An interpolation scheme has \emph{interpolation locality}\index{interpolation
locality} if a change in one input quote changes the curve only near that
input's maturity. Locality decides what a bucketed hedge looks like: under a
local scheme an \omterm{def:rc:curve-construction:pillar}{off-pillar} swap is hedged with the \omterm{def:rc:curve-construction:pillar}{pillars} around it; under a
non-local one it acquires risk, of either sign, on \omterm{def:rc:curve-construction:pillar}{pillars} far from its
maturity.
\end{definition}

\begin{omfigure}
\begin{tikzpicture}
\begin{axis}[width=12cm,height=5.8cm,
  xlabel={maturity $T$ (years)},ylabel={change in $f(0,T)$ (bp)},
  tick label style={font=\small},label style={font=\small},
  xmin=0,xmax=20,ymin=-5.5,ymax=5.5,grid=major,grid style={gray!25},
  legend columns=4,legend style={font=\footnotesize,at={(0.5,-0.25)},anchor=north,draw=none,fill=none,/tikz/every even column/.append style={column sep=8pt}},
  table/col sep=comma]
\addplot[gray,thick] table[x=t,y=lz]{figdata/rates-credit-risk/01-curve-construction/locality.csv};
\addplot[omDef,thick] table[x=t,y=ff]{figdata/rates-credit-risk/01-curve-construction/locality.csv};
\addplot[omProp,thick,dashed] table[x=t,y=cs]{figdata/rates-credit-risk/01-curve-construction/locality.csv};
\addplot[omThm,very thick] table[x=t,y=mc]{figdata/rates-credit-risk/01-curve-construction/locality.csv};
\legend{linear zero,flat forward,cubic on zeros,monotone convex}
\end{axis}
\end{tikzpicture}

\omcaption{The response of the forward curve to a one-basis-point rise in the
seven-year swap quote, with the curve recalibrated. Flat forwards move only
between five and ten years (up before seven, down after, so that the ten-year
swap still reprices); the monotone convex curve adds small
ripples just outside that span; the cubic spline moves the forwards from three
years to beyond twenty. Data: the chapter's tutorial.}\label{fig:rc:curve-construction:locality}
\end{omfigure}

\begin{example}[An eight-year swap on four curves]\label{ex:rc:curve-construction:eight}
A par payer swap of USD~100 million for eight years has a parallel DV01 of
about USD~69\,600 on all four curves: they agree on the level. They disagree
on where it sits (\cref{fig:rc:curve-construction:buckets}). Flat forwards put
USD~40\,400 on the seven-year quote and USD~29\,200 on the ten-year; the cubic
spline puts USD~62\,100 on seven years, $-22\,200$ on five, $+11\,000$ on four
and $-9\,500$ on twelve; monotone convex lies in between. Each set of buckets
sums to the parallel DV01, and each asks for a different hedge.
\end{example}

\begin{omfigure}
\begin{tikzpicture}
\begin{axis}[width=12cm,height=5.4cm,ybar,bar width=7pt,
  ylabel={DV01 (USD thousand per bp)},
  tick label style={font=\small},label style={font=\small},
  xmin=-0.6,xmax=6.6,ymin=-30,ymax=70,grid=major,grid style={gray!25},
  xtick=data,xticklabels from table={figdata/rates-credit-risk/01-curve-construction/buckets8.csv}{tenor},
  legend columns=3,legend style={font=\footnotesize,at={(0.5,-0.2)},anchor=north,draw=none,fill=none,/tikz/every even column/.append style={column sep=8pt}},
  table/col sep=comma]
\addplot[fill=omDef!55,draw=omDef] table[x=i,y=ff]{figdata/rates-credit-risk/01-curve-construction/buckets8.csv};
\addplot[fill=omProp!55,draw=omProp] table[x=i,y=cs]{figdata/rates-credit-risk/01-curve-construction/buckets8.csv};
\addplot[fill=omThm!55,draw=omThm] table[x=i,y=mc]{figdata/rates-credit-risk/01-curve-construction/buckets8.csv};
\legend{flat forward,cubic on zeros,monotone convex}
\end{axis}
\end{tikzpicture}

\omcaption{Bucketed DV01 of an eight-year par payer swap of USD~100 million by
input quote, for three interpolations (each quote bumped by a hundredth of a
basis point, curve recalibrated, result scaled to one basis point). The
eight-year date lies between the seven- and ten-year \omterm{def:rc:curve-construction:pillar}{pillars}; flat forwards
load only those two, monotone convex adds small five- and twelve-year buckets,
the cubic spline large ones of both signs from three to fifteen years. Data:
the chapter's tutorial.}\label{fig:rc:curve-construction:buckets}
\end{omfigure}

\begin{remark}[Monotone convex is not linear in its inputs]\label{rem:rc:curve-construction:nonlinear}
Linear-zero, flat-forward and spline curves are linear functions of the \omterm{def:rc:curve-construction:pillar}{pillar}
rates, so a bucketed DV01 does not depend on the size of the bump. The
monotone convex curve is not: which of its four cases applies on an interval
depends on the inputs, so a one-basis-point bump can switch a case and move
the forward shape inside an interval by more than the bump itself. The value
stays continuous, but the buckets change with the bump size. The chapter
therefore bumps by a hundredth of a basis point and scales; exercise~7
measures the difference.
\end{remark}

\section{Tutorial: one set of quotes, four curves}\label{tut:rc:curve-construction:tutorial}

\begin{tutorial}
\textbf{Goal.} Calibrate one dollar curve from deposits, futures and swaps
under four interpolations, and compare forwards, locality and hedges.
\textbf{End state:}
\cref{fig:rc:curve-construction:forwards,fig:rc:curve-construction:locality,fig:rc:curve-construction:buckets}
and the numbers of \cref{ex:rc:curve-construction:eight}.
\begin{tutsteps}
\item \textbf{Instruments.} Two deposits, four futures from the third
Wednesday of December 2026, eleven swaps; no one-year swap.
\omcode{code/rates-credit-risk/01-curve-construction/python/rc_curves.py}{30}{36}{The instrument set of the chapter's curve.}\label{lst:rc:curve-construction:instruments}
\item \textbf{Calibrate} with \texttt{calibrate(SPOT, instruments(), kind)} for each
of the four kinds; check that the maximum residual is below $10^{-12}$ and the
number of iterations is three.
\item \textbf{Forwards.} Tabulate \texttt{curve.fwd\_t(t)} every 0.05 years; you
should see the sawtooth of linear zeros at the two-year \omterm{def:rc:curve-construction:pillar}{pillar} (2.964\% then
3.331\%).
\item \textbf{Locality and buckets.} Run \texttt{locality} for the \omterm{def:rc:curve-construction:pillar}{pillar} \texttt{7Y} and
\texttt{off\_pillar\_buckets(8)}; \texttt{fig\_rc\_curves.py} writes the charts.
\end{tutsteps}
\textbf{What to change next.} Put the one-year swap back two basis points
away from the futures and look at the forward curve in September 2027; bump by a full basis point
under monotone convex and compare the buckets with
\cref{fig:rc:curve-construction:buckets}.
\end{tutorial}

\section{Build: a multi-instrument curve builder}\label{bld:rc:curve-construction:curvebuild}

\begin{build}
\bfield{Purpose} The first curve of Book 6's library: every rates instrument,
exposure simulation and risk run of the book reads a curve built here, and the
pricing library of One Quant Book 5, chapter 28, accepts it as a discount
curve.
\bfield{Interface} \texttt{Deposit}, \texttt{Future}, \texttt{Swap} with
\texttt{quote()} and \texttt{model(curve)}; \texttt{calibrate(spot, instruments,
kind)} returning the curve, the Jacobian and diagnostics; \texttt{ZeroCurve}
with \texttt{df}, \texttt{zero\_t}, \texttt{fwd\_t} and \texttt{bumped(pillar,
size)}; \texttt{bucket\_sensitivities}; kinds \texttt{linear\_zero},
\texttt{flat\_forward}, \texttt{cubic\_zero}, \texttt{monotone\_convex}.
\bfield{Rules} \omterm{def:rc:curve-construction:pillar}{Pillars} at instrument maturities, strictly increasing (an
overlap is an error); ACT/365 times, ACT/360 accruals; flat forward beyond the
last \omterm{def:rc:curve-construction:pillar}{pillar}; Book 2's \texttt{firm\_curve} imported, not edited.
\bfield{Acceptance tests} \texttt{code/firm/curvebuild/tests/}: every
instrument reprices under every kind; flat-forward on swaps equals Book 2's
bootstrap; flat par quotes give flat forwards; the monotone convex forward is
continuous at \omterm{def:rc:curve-construction:pillar}{pillars}; local kinds have a lower-triangular Jacobian and the
spline does not; a seven-year bump spreads beyond fifteen years only under the
spline; \texttt{bumped} shifts one \omterm{def:rc:curve-construction:pillar}{pillar} or all.
\bfield{Stretch} Turn and meeting-date \omterm{def:rc:curve-construction:pillar}{pillars}; a least-squares mode with more
instruments than \omterm{def:rc:curve-construction:pillar}{pillars}; the positivity amendment of the monotone convex
method; an analytic Jacobian.
\end{build}

\begin{omsources}
\item P.\ S.\ Hagan and G.\ West, ``Interpolation methods for curve
construction'', \emph{Applied Mathematical Finance} 13(2), 2006.
\item C.\ R.\ Nelson and A.\ F.\ Siegel, ``Parsimonious modeling of yield
curves'', \emph{Journal of Business} 60(4), 1987; L.\ E.\ O.\ Svensson,
``Estimating and interpreting forward interest rates: Sweden 1992--1994'', NBER
Working Paper 4871, 1994.
\item US Department of the Treasury, \emph{Treasury Yield Curve Methodology};
ECB, \emph{Euro area yield curves} (ECB Data Portal).
\item L.\ Andersen and V.\ Piterbarg, \emph{Interest Rate Modeling}, volume 1,
Atlantic Financial Press, 2010, chapter 6.
\end{omsources}

\section{Exercises}

\begin{exercise}[$\star$]\label{exo:rc:curve-construction:1}
The continuously compounded zero rates at two and three years are $3.25\%$
and $3.30\%$. Under \omterm{def:rc:curve-construction:flatfwd}{flat-forward interpolation}, what is the \omterm{def:rc:curve-construction:pillar}{instantaneous
forward rate} at $2.5$ years?
\end{exercise}

\begin{exercise}[$\star$]\label{exo:rc:curve-construction:2}
A SOFR future on the quarter from June to September 2027 trades at $96.72$ and
its convexity adjustment is $0.6$ basis points. What forward rate must the curve
reproduce for that quarter?
\end{exercise}

\begin{exercise}[$\star$]\label{exo:rc:curve-construction:3}
Why does the chapter's curve leave out the one-year swap, and what would
happen to the calibration if it were put back with a quote two basis points
away from what the futures imply?
\end{exercise}

\begin{exercise}[$\star\star$]\label{exo:rc:curve-construction:4}
On the chapter's linear-zero curve the \omterm{def:rc:curve-construction:pillar}{pillars} around two years are at
$(1.2110, 3.4165\%)$, $(2.0027, 3.2883\%)$ and $(3.0027, 3.3095\%)$ (years,
zero rate). Compute the instantaneous forward just before and just after the
two-year \omterm{def:rc:curve-construction:pillar}{pillar}.
\end{exercise}

\begin{exercise}[$\star\star$]\label{exo:rc:curve-construction:5}
From the ECB parameters of \cref{ex:rc:curve-construction:ecb}, compute the
ten-year \omterm{def:rc:curve-construction:pillar}{instantaneous forward rate} and the limit of the zero rate as $T\to 0$.
\end{exercise}

\begin{exercise}[$\star\star$]\label{exo:rc:curve-construction:6}
Read \cref{fig:rc:curve-construction:buckets}. A trader hedges the eight-year
swap with the spline's buckets, then the desk moves to flat forwards. Which
hedges become unnecessary, and why did the spline ask for them?
\end{exercise}

\begin{exercise}[$\star\star\star$]\label{exo:rc:curve-construction:7}
\emph{Coding.} With \texttt{off\_pillar\_buckets(8, bump=1e-4)}, compute the
monotone convex buckets of the eight-year swap with a full-basis-point bump and
their sum; compare with the sum for a hundredth-of-a-basis-point bump and with
the parallel DV01, and explain.
\end{exercise}

\begin{exercise}[$\star\star\star$]\label{exo:rc:curve-construction:8}
\emph{Find the flaw.} ``Our new curve reprices every input to $10^{-12}$, so
its prices for every other instrument are right.'' Correct it.
\end{exercise}

\section{Problem: The Monday Loss}

\begin{problem}[{Weekend problem --- an interpolation changed over a weekend}]\label{pb:rc:curve-construction:1}
A desk's book holds seven par swaps traded on the chapter's curve under
linear-zero interpolation: payers of 6, 9, 13 and 22 years (USD~0.9, 0.6, 0.5
and 0.8 billion) and receivers of 8, 11 and 17 years (USD~0.7, 0.8 and 0.7
billion), USD~5.0 billion in all. Over the weekend the curve is rebuilt from
the same quotes with \omterm{def:rc:curve-construction:monotone}{monotone convex interpolation}.

\medskip\noindent\textbf{Part I --- The change.}
\begin{enumerate}
\item Why is each swap worth exactly zero on Friday?
\item Does any input instrument change value on Monday? Why not?
\item The 17-year par rate moves from $3.9975\%$ to $4.0074\%$. What does that
say about the forwards between 15 and 20 years?
\item Which swaps of the book are exposed to the change, and which are not?
\item Give the six-year par rate before and after, and explain why it barely
moves.
\end{enumerate}

\medskip\noindent\textbf{Part II --- The P\&L.}
\begin{enumerate}[resume]
\item Give the book's value on Monday.
\item Is it a profit or a loss, and which swaps produce most of it?
\item Should the desk book it as trading P\&L? Who should decide?
\item How would you present it to the head of desk in one sentence?
\item Would a book of pillar-maturity swaps have shown anything?
\end{enumerate}

\medskip\noindent\textbf{Part III --- The hedges.}
\begin{enumerate}[resume]
\item The fifteen-year bucket was $-\text{USD}~346\,584$ per basis point and is now
$-583\,281$. The fifteen-year par swap has a DV01 of USD~1\,149\,075 per billion.
Give the hedge notional before and after.
\item Do the same for the twenty-year bucket ($345\,841$ then $465\,177$; USD~1\,388\,775
per billion).
\item The sum of the buckets is USD~555\,270 before and 557\,510 after. Why so
close, when single buckets moved by a third or more?
\item Why were the buckets computed with a hundredth-of-a-basis-point bump?
\item A new twelve-year bucket appears ($-77\,741$). Where did it come from?
\end{enumerate}

\medskip\noindent\textbf{Part IV --- Judgement.}
\begin{enumerate}[resume]
\item Give two arguments for the new interpolation and one against.
\item How should a change of curve methodology be governed?
\item What reserve might a controller hold for curve-methodology uncertainty?
\item State the \emph{named result}: the P\&L of the switch and the change in the
fifteen-year hedge.
\item In one sentence: what does a curve's interpolation decide?
\end{enumerate}
\end{problem}

\section{Interview questions}

\begin{interviewq}[$\star$ \iqroles{researcher, bank}]\label{iq:rc:curve-construction:1}
What is the difference between curve bootstrapping and a global curve fit, and
when do you need the second?
\end{interviewq}

\begin{interviewq}[$\star$ \iqroles{trader, researcher}]\label{iq:rc:curve-construction:2}
Why does linear interpolation of zero rates give a bad forward curve?
\end{interviewq}

\begin{interviewq}[$\star\star$ \iqroles{researcher, developer}]\label{iq:rc:curve-construction:3}
List the properties you would want from a curve interpolation, and say which
ones conflict.
\end{interviewq}

\begin{interviewq}[$\star\star$ \iqroles{trader}]\label{iq:rc:curve-construction:4}
Your eight-year swap shows risk on the four-year and twelve-year buckets. What
do you suspect?
\end{interviewq}

\begin{interviewq}[$\star\star$ \iqroles{developer, risk}]\label{iq:rc:curve-construction:5}
How do you handle the year-end turn and central-bank meeting dates in a curve
builder?
\end{interviewq}

\begin{interviewq}[$\star\star\star$ \iqroles{researcher, risk}]\label{iq:rc:curve-construction:6}
Your bucketed deltas change when you change the bump size from one basis point
to a tenth. Is the code wrong?
\end{interviewq}
